\documentclass{article}
\usepackage[T1]{fontenc}
\usepackage{lmodern}
\usepackage{graphicx}
\usepackage{amsmath,amssymb}
\usepackage[a4paper,margin=2cm]{geometry}
\usepackage[absolute,overlay]{textpos}
\setlength{\TPHorizModule}{1mm}
\setlength{\TPVertModule}{1mm}
\usepackage[indonesian]{babel}
\usepackage{hyperref}
\setlength{\emergencystretch}{2em}
% unit: Halaman 85

\begin{document}

% === Halaman 85 (llm) ===

\textbf{Step 5:} State hasil 10 putaran-ganda dijumlahkan dengan state awal yang disimpan pada step 3:

\[
\mathbf{y = (x' + x) \pmod{2^{32}}}
\]

Jika plainteks berukuran n x 512 bit, maka Salsa20 dieksekusi sebanyak n kali, dengan \textit{counter} bertambah 1, sedangkan \textit{nonce} tetap.

\textbf{Kode program C++:}

\begin{verbatim}
#include <stdint.h>
#define ROTL(a,b) (((a) << ((b)) | ((a) >> (32 - (b)))))
#define QR(a, b, c, d)(\ 
    b ^= ROTL(a + d, 7), \ 
    c ^= ROTL(b + a, 9), \ 
    d ^= ROTL(c + b, 13), \ 
    a ^= ROTL(d + c, 18))
#define ROUNDS 20
\end{verbatim}

\end{document}
